% Nouveau EGA V (cote 25): archive cote 25, site page 164 (PDF page 165).
% typescript. Draft compared: none.
% AI-transcribed by Claude Opus 5.5 (claude-opus-5-5, adaptive thinking, effort high), Message Batches API, batch ega25b (msgbatch_01A7n5YuXubUPNK1zsepi8RG), 2026-11-13; not yet checked by eye.
\scanpage{25/164}
\mspage{7}

k-isomorphe à $k[T]/\unsure{(}T^n)$ . Si $d=1$ ,\inserted[typed]{x rat. sur k,} la notion de point singulier
\xout{ordinaire} élémentaire de multiplicité n correspond à la terminologie
classique \xout{\illegible{1}}``point singulier ordinaire à n tangente\unsure{\xout{s}|s} disctinctes'' .
\margin{(}(On aura \inserted{déjà} explicité dans 15.2. que si x est un point de X rat sur k, alors
la notion de singularité élémentaire de mult. n, au sens absolu ou au
sens relatif, \xout{\unsure{elle}} est la même; cela reste d'ailleurs valable si k(x)
est une extension finie séparable de k - sans doute d'ailleurs ``finie''
est inutile - et ce fait mérite d'être inséré en corollaire \xout{\unsure{si on arrive}}.
ou en proposition\hand{)}.

NB J'avais inclus au par.20 un nº sur l'éclatement d'un
préschéma X le long d'un sous-préschéma \inserted{fermé} Y tel que l'augmentation
$\underline{O}_X \longrightarrow \underline{O}_Y$ soit élémentaire de mult.n (auquel j'ai fait allusion dans
15.6.). Je signale qu'il s'agit d'un nº court et bien venu, dont les
seuls ingrédients sont les résultats généraux de lissité du par.17,et
la définition \inserted{du début de} 17.15.\unsure{2}. \inserted[typed]{(et encore pourrait-on éviter cette dernière).} Rien n'empêcherait d'incopporer ces résultats
dès ici, \xout{\unsure{par dans}} par exemple immédiatement après la définition
\xout{17}15.1. ci-dessus; alors \inserted[typed]{(la gamme ``géométrique''de 15.2. et la suite} 15.\unsure{3}. \unsure{à} 15.7. pourraient être séparés
de ces \xout{\unsure{calculs}} résultats, en les groupant dans un Nº 17.16, \xout{\illegible{2}}
\xout{\unsure{lorsque}} De toutes façons, je m'aperçois en cours de rédaction que le
\margin{(}nouveau nº 15 prévu sur les formes lisses éclate en \inserted{au moins} deux numéros assez
distincts et indépendants, l'un contenant des sorites généraux pour
les trucs ``élémentaires de mult n'' , n quelconque, l'autre \deleted{qui} cont\over{ient}{enant} la
caractérisation des formes \uline{quadratiques} lisses, et qui n'emprunte au
précédent que le nº 15.1. , i.e. pratiquement la définition de la lissité
d'une forme.

\mspage{164}
